% Generated draft: requires mathematical and source review.
\documentclass[11pt]{article}
\usepackage[margin=25mm]{geometry}
\usepackage{fontspec}
\setmainfont{Latin Modern Roman}
\usepackage{amsmath,amssymb,mathtools,unicode-math}
\setmathfont{Latin Modern Math}
\usepackage{enumitem,needspace,xurl,hyperref}
\setlength{\parindent}{0pt}
\setlength{\parskip}{5pt}
\setlength{\emergencystretch}{3em}
\newcommand{\N}{\mathbb{N}}
\newcommand{\Z}{\mathbb{Z}}
\newcommand{\Q}{\mathbb{Q}}
\newcommand{\R}{\mathbb{R}}
\newcommand{\C}{\mathbb{C}}
\newcommand{\F}{\mathbb{F}}
\newcommand{\A}{\mathbb{A}}
\newcommand{\PP}{\mathbb{P}}
\newcommand{\E}{\mathbb{E}}
\newcommand{\eps}{\varepsilon}
\newcommand{\dd}{\,\mathrm{d}}
\begin{document}
\section*{UnsolvedMath ID: 800033}
\textbf{Source problem code:} catalogue entry
\subsection{Definitions and mathematical statement}
A semigroup \(S\) is \emph{regular} if
\[
\forall s\in S\ \exists t\in S:\quad sts=s.
\]
A group \(G\) is \emph{periodic} if every element has finite order:
\[
\forall g\in G\ \exists n\in\Z_{\geq 1}:\quad g^n=1_G.
\]
Under the interpretation of the source's phrase ``power one'' as the power semigroup of nonempty subsets, put
\[
\mathcal P_+(G)=\{A\subseteq G:A\ne\varnothing\},
\qquad AB=\{ab:a\in A,\ b\in B\}.
\]
An embedding \(\phi:S\hookrightarrow\mathcal P_+(G)\) is an injective map satisfying \(\phi(st)=\phi(s)\phi(t)\) for all \(s,t\in S\).

The requested characterization, conditional on this interpretation, is: find necessary and sufficient conditions on a regular semigroup \(S\) for
\[
\exists\text{ a periodic group }G\ \exists\text{ an injective semigroup homomorphism }
\phi:S\longrightarrow\mathcal P_+(G).
\]
The supplied excerpt does not define ``power one'' or specify whether the empty subset is admitted. If the full power set is intended, replace \(\mathcal P_+(G)\) by \(\mathcal P(G)\), with the same setwise multiplication. These conventions require source confirmation; no equivalence of the resulting embedding conditions is asserted here.

A semigroup is \emph{completely regular} if every element belongs to a subgroup of the semigroup. The source reports that, for completely regular semigroups, the necessary and sufficient conditions are ``inversion and periodicity,'' citing \([\mathrm{Bersh+Kubl}]\). The excerpt does not define ``inversion.'' One possible interpretation is the inverse-semigroup property: for every \(s\in S\), there exists a unique \(t\in S\) such that \(sts=s\) and \(tst=t\). This interpretation remains unconfirmed. In the usual semigroup terminology, periodicity means that \(\{s^n:n\geq1\}\) is finite for every \(s\in S\). The reported completely regular case is retained as supplied background, with its terminology and reference unresolved.
\subsection{Short English statement}
Characterize the regular semigroups that embed into a power semigroup of a periodic group. The subset convention and the reported completely regular case need clarification, especially the meaning of ``inversion.''
\subsection{Sources}
\begin{itemize}
\item Supplied local UnsolvedMath registry record: permanent ID \(800033\), source code \texttt{catalogue entry}, title \emph{The Lyapin's Notebook---Problem 35}. Its status is supplied as \texttt{open}; this is not an independently verified status assessment.
\item Bershadsky, Kublanovsky, and Mashevitzky, \emph{The Lyapin's Notebook: Unsolved Problems in Semigroup Theory} (2026), Problem 35. Supplied source URL: \url{https://arxiv.org/abs/2604.04763}. The record reports access on 2026-10-16 and extraction by \texttt{arxiv\_latex}; neither access nor source contents were independently verified here.
\item The statement attributes the problem to S.G. Bershadsky and S.I. Kublanovsky. The registry separately lists Bershadsky, Kublanovsky, and Mashevitzky under \texttt{proposed\_by} and gives \(2026\) under \texttt{proposed\_year}; these metadata do not establish the original proposal date or resolve the attribution difference.
\item The supplied background cites \([\mathrm{Bersh+Kubl}]\) for the completely regular case, but supplies no bibliographic details identifying that reference.
\end{itemize}
\end{document}
